% ── SECTION 6: THE SEMANTIC BRIDGE (~600 words) ──────────────────────────────

\section{The Semantic Bridge: One Phenomenon, Three Names}
\label{sec:semantic}

The most important single insight of this project is this:
\mahabbat is not a human emotion projected onto an indifferent universe.
It is the fundamental property of reality --- the force that forms and sustains
relationship at every scale --- which human beings experience, in its most
conscious and intense form, as emotion.

Table~\ref{tab:semantic} presents the semantic bridge:
each tradition's technical term for the same phenomenon, with the key source
and the thesis it supports.

\begin{table}[h]
\centering
\caption{The semantic bridge: one phenomenon, five names.}
\label{tab:semantic}
\begin{tabularx}{\textwidth}{lXXl}
\toprule
\textbf{Scale / Domain} & \textbf{Technical Term} & \textbf{What It Names} & \textbf{Thesis} \\
\midrule
Quantum (micro)     & Entanglement / Affinity Tensor $\alpha$
                    & Nonlocal relational bond; measure of coherence & T1 \\
Cosmic (macro)      & Gravity / Spacetime geometry
                    & Universally attractive force; Einstein's geometry of relation & T1 \\
Process Philosophy  & Prehension / Eros
                    & Relational act of becoming; creative drive toward richer relation & T2 \\
Bahá'í Writings     & \mahabbat / Mawaddat / Affinity
                    & Creative foundation of all existence at every degree & T3 \\
Human experience    & Love
                    & Most conscious form of the universal relational force & All \\
\bottomrule
\end{tabularx}
\end{table}

Note what is \emph{not} being claimed.
We are not asking physicists to call entanglement ``love,'' or Bahá'í scholars
to redefine \mahabbat as the Affinity Tensor.
Each tradition keeps its own vocabulary, its own methodology, and its own
epistemic standards.
What we claim is that when each tradition's instruments point at the same
structure, from independent directions, across a 170-year sequence,
that structure is real.

The table has a deliberate architecture.
Reading vertically, the five rows trace the same phenomenon across five
domains of inquiry --- from the subatomic to the cosmic, from the formal
to the experiential.
Reading horizontally, each row names the same relational force in the
language native to its domain.
No row is more fundamental than the others; the Affinity Tensor is not
a more precise version of \mahabbat, any more than gravity is a more
precise version of eros.
Each is the correct technical term for the relational force at the
scale and in the vocabulary where that tradition operates.

The deepest implication of the semantic bridge is not terminological
but ontological.
If five traditions, operating at five different scales and by five different
methods, all name the same force --- if physicists call it entanglement
at the quantum scale, Einstein called it spacetime geometry at the cosmic
scale, Whitehead called it eros in the philosophical vocabulary of process,
the Bahá'í writings call it \mahabbat in the vocabulary of revelation, and
human beings call it love in the vocabulary of experience --- then the most
parsimonious explanation is that the force is real and singular, and that
the diversity of names reflects the diversity of the instruments trained
upon it, not a diversity of things.

This is not a new claim in the history of philosophy.
The intuition that love is a cosmic and not merely a personal force is
ancient.
What this project adds is precision: a formal measure (the Affinity Tensor
$\alpha$), a testable prediction (the $\URA$ gate), and a three-tradition
independent confirmation that triangulates the target from different
directions simultaneously.
The argument is not that love is important.
The argument is that love --- understood as the force of relational affinity
that constitutes and sustains existence at every scale --- is what reality
is made of.
That argument is now formal enough to be tested, and it has been
independently arrived at by three traditions that did not communicate
with each other.
